\documentclass[12pt,letterpaper]{ctexart}

% ===== 数学建模论文模板 =====
% 基于实际竞赛论文格式解析
% 适用于国赛、美赛(MCM/ICM)、华中杯等

% ---------- 页面设置 ----------
% 页边距：左3.17cm 右2.75cm 上下约2.6cm
\usepackage[top=2.6cm,bottom=2.2cm,left=3.17cm,right=2.75cm]{geometry}

% ---------- 数学与公式 ----------
\usepackage{amsmath,amssymb,amsthm}

% ---------- 图表 ----------
\usepackage{graphicx,subfig}
\usepackage{booktabs,multirow,array,longtable}
\usepackage{float}

% ---------- 算法伪代码 ----------
\usepackage{algorithm,algorithmic}
\renewcommand{\algorithmicrequire}{\textbf{输入:}}
\renewcommand{\algorithmicensure}{\textbf{输出:}}

% ---------- 引用与链接 ----------
\usepackage{hyperref}
\hypersetup{
    colorlinks=true,
    linkcolor=black,
    citecolor=black,
    urlcolor=black
}

% ---------- 列表 ----------
\usepackage{enumitem}

% ---------- 图表标题 ----------
\usepackage{caption}
\captionsetup{
    font=small,
    labelfont=bf,
    labelsep=quad,
    justification=centering
}

% ---------- 页眉页脚 ----------
\usepackage{fancyhdr}
\pagestyle{fancy}
\fancyhf{}
\fancyhead[C]{\songti\zihao{5}\leftmark}     % 页眉：居中，宋体五号，显示当前节名
\fancyfoot[C]{\thepage}                        % 页脚：居中页码
\renewcommand{\headrulewidth}{0.4pt}           % 页眉线宽度
\setlength{\headheight}{14pt}

% ---------- 首行缩进 ----------
\usepackage{indentfirst}
\setlength{\parindent}{2em}                    % 首行缩进2个汉字

% ---------- 行距 ----------
\linespread{1.3}                               % 1.3倍行距

% ---------- 标题格式设置 ----------
% 各级标题使用 ctex 默认，通过 \zihao 控制字号：
% 三号=16pt, 四号=14pt, 小四=12pt

% 一级标题：四号(14pt)黑体居中
\CTEXsetup[name={,、},number={\chinese{section}},format={\heiti\zihao{4}\centering},aftername={}]{section}

% 二级标题：小四(12pt)黑体左对齐
\CTEXsetup[name={,},number={\arabic{subsection}},format={\heiti\zihao{-4}\raggedright},aftername={}]{subsection}

% 三级标题：小四(12pt)黑体左对齐
\CTEXsetup[name={,},number={\arabic{subsection}.\arabic{subsubsection}},format={\heiti\zihao{-4}\raggedright},aftername={}]{subsubsection}

% 图表编号：按节编号（如 表5-1、图5-1）
\numberwithin{equation}{section}
\numberwithin{figure}{section}
\numberwithin{table}{section}

% 重定义图表标题标签为中文
\renewcommand{\figurename}{图}
\renewcommand{\tablename}{表}
\renewcommand{\algorithmcfname}{算法}

\begin{document}

% ==========================================
% 首页：标题 + 摘要 + 关键词
% ==========================================

% --- 标题：16pt(三号)黑体加粗居中 ---
\begin{center}
{\heiti\zihao{3}\textbf{【在此处填写论文标题】}}

\vspace{1cm}

% --- 摘要标签：16pt(三号)黑体居中 ---
{\heiti\zihao{3}摘\ \ 要}
\end{center}

\vspace{0.5cm}

% --- 摘要正文：14pt(四号)宋体，首行缩进2字符 ---
% 摘要要求：300-500字，包含问题背景 + 方法 + 结果 + 创新点
{
\songti\zihao{4}
【在此处填写摘要正文。摘要应包含以下四要素：
(1) 问题背景：一句话说明问题的实际意义；
(2) 主要方法：简述本文采用的建模方法和算法；
(3) 关键结果：给出主要的数值结果和发现；
(4) 创新点：本文工作的贡献或创新之处。】

\vspace{1cm}

% --- 关键词：14pt(四号)，"关键词："加粗 ---
\noindent{\textbf{关键词：}}【关键词1；关键词2；关键词3；关键词4；关键词5】
}

\newpage

% ==========================================
% 目录页
% ==========================================
{
\pagestyle{empty}
\tableofcontents
}
\newpage

% ==========================================
% 正文开始，重置页码
% ==========================================
\setcounter{page}{1}
\pagestyle{fancy}

% ==========================================
% 一、问题重述
% ==========================================
\section{问题重述}

\subsection{问题背景}
【在此处描述问题的实际背景和意义。结合行业发展趋势、政策背景等进行阐述，为后续建模提供现实依据。约300-400字。】

\subsection{问题重述}
【在此处用学术语言准确重述题目给出的问题。逐问说明已知条件、求解目标和约束限制。注意不要照抄题目原文，应经过自己的理解重新表述。约400-500字。】

% ==========================================
% 二、问题分析
% ==========================================
\section{问题分析}

\subsection{问题1分析}
【分析问题1的核心难点、关键约束，指出解题的关键思路。说明：(1) 问题类型和复杂性 (如NP-hard)；(2) 主要挑战点；(3) 解题思路框架。约200-300字。】

\subsection{问题2分析}
【分析问题2与问题1的差异，新增约束带来的影响，以及解题策略的调整。约200-300字。】

\subsection{问题3分析}
【分析问题3的动态性特点，实时性要求和局部性策略。约200-300字。】

% ==========================================
% 三、模型假设
% ==========================================
\section{模型假设}
【列出模型的假设条件，编号说明。每条假设一句话清晰表述。约5-8条假设。】

\begin{enumerate}[label=\arabic*.]
    \item 【假设1：如忽略次要因素，将问题进行合理简化】
    \item 【假设2：如假设某参数为常数，不随时间变化】
    \item 【假设3：如不考虑极端天气等特殊情况的影响】
    \item 【假设4】
    \item 【假设5】
\end{enumerate}

% ==========================================
% 四、参数与变量说明
% ==========================================
\section{参数与变量说明}

【将模型中使用的所有参数和决策变量以表格形式列出。参数表和变量表分列。】

% --- 参数表 ---
\begin{table}[H]
\centering
\caption{模型参数}
\begin{tabular}{c|l|c}
\toprule
\textbf{符号} & \textbf{含义} & \textbf{单位} \\
\midrule
$N$ & 【含义】 & 【单位】 \\
$c_{ij}$ & 【含义】 & 【单位】 \\
\bottomrule
\end{tabular}
\end{table}

% --- 决策变量表 ---
\begin{table}[H]
\centering
\caption{决策变量}
\begin{tabular}{c|l|c}
\toprule
\textbf{符号} & \textbf{含义} & \textbf{类型} \\
\midrule
$x_{ijk}$ & 【含义】 & 0-1变量 \\
$s_{ik}$ & 【含义】 & 连续变量 \\
\bottomrule
\end{tabular}
\end{table}

% ==========================================
% 五、模型的建立与求解（核心章节）
% ==========================================
\section{模型的建立与求解}

% --- 5.1 数据预处理 ---
\subsection{数据预处理}

\subsubsection{客户分布与需求特征}
【在此描述数据来源、数据概况、预处理过程（订单聚合、异常处理、特征分析等）。用表格和图展示关键统计结果。约400-600字。】

% 示例：插入图片
%\begin{figure}[H]
%    \centering
%    \includegraphics[width=0.75\textwidth]{figure1_customer_distribution.png}
%    \caption{客户空间分布与绿色配送区}
%    \label{fig:customer-dist}
%\end{figure}

\subsubsection{速度时段建模}
【如有时间相关变量，说明时段划分方法和速度模型。约200-300字。】

\subsubsection{能耗建模与负载修正}
【说明能耗计算模型，包括负载修正系数、跨时段积分方法等。约200-300字。】

% --- 5.2 问题1的模型与求解 ---
\subsection{问题1：静态场景下的车辆调度}

\subsubsection{模型建立}

【说明建模思路】考虑【问题背景】，建立如下数学模型：

\paragraph{目标函数}

\begin{equation}\label{eq:obj}
\min Z = C_{\text{vehicle}} + C_{\text{energy}} + C_{\text{penalty}} + C_{\text{carbon}}
\end{equation}

其中：
\begin{itemize}
    \item $C_{\text{vehicle}}$ —— 【车辆启动成本，计算公式为 $C_{\text{vehicle}} = \sum_{k \in K} c_k^{\text{fixed}} \cdot y_{0k}$】；
    \item $C_{\text{energy}}$ —— 【能耗成本，包括燃油费和电费】；
    \item $C_{\text{penalty}}$ —— 【时间窗违反惩罚成本】；
    \item $C_{\text{carbon}}$ —— 【碳排放成本（碳税）】。
\end{itemize}

\paragraph{约束条件}

\begin{enumerate}[label=(\arabic*)]
    \item \textbf{容量约束（重量和体积）}：
    \begin{align}
        \sum_{i \in N} w_i \cdot y_{ik} &\leq W_k, \quad \forall k \in K \label{eq:cap-w}\\
        \sum_{i \in N} v_i \cdot y_{ik} &\leq V_k, \quad \forall k \in K \label{eq:cap-v}
    \end{align}
    约束(\ref{eq:cap-w})确保每辆车的载重量不超过其额定载重，约束(\ref{eq:cap-v})确保装载体积不超过车辆容积。这两条约束是车辆路径问题的基本约束。

    \item \textbf{服务完整性约束}：
    \begin{equation}
        \sum_{k \in K} y_{ik} = 1, \quad \forall i \in N \label{eq:service}
    \end{equation}
    约束(\ref{eq:service})确保每个客户节点恰好被一辆车服务一次，不允许遗漏或重复服务。

    \item \textbf{路径连续性约束}：
    \begin{align}
        \sum_{j \in N \cup \{0\}} x_{ijk} &= y_{ik}, \quad \forall i \in N, k \in K \label{eq:flow1}\\
        \sum_{i \in N \cup \{0\}} x_{ijk} &= y_{jk}, \quad \forall j \in N, k \in K \label{eq:flow2}
    \end{align}
    约束(\ref{eq:flow1})和(\ref{eq:flow2})确保每条路径是连续的——车辆进入某节点后必须从该节点离开。

    \item \textbf{配送中心约束}：
    \begin{equation}
        \sum_{j \in N} x_{0jk} = \sum_{i \in N} x_{i0k} = 1, \quad \forall k \in K \label{eq:depot}
    \end{equation}
    约束(\ref{eq:depot})要求每辆车从配送中心出发，完成配送后返回配送中心。

    \item \textbf{时间窗软约束}：
    \begin{equation}
        a_i \leq s_{ik} \leq b_i + \delta_{\text{late}}, \quad \forall i \in N, k \in K \label{eq:tw}
    \end{equation}
    约束(\ref{eq:tw})为时间窗软约束，允许延迟到达但产生惩罚成本。其中 $a_i$ 和 $b_i$ 分别为客户 $i$ 的最早和最晚服务时间，$\delta_{\text{late}}$ 为允许的最大延迟。早到等待成本为 $c_{\text{early}}=20$元/h，晚到惩罚成本为 $c_{\text{late}}=50$元/h。

    \item \textbf{车队数量约束}：
    \begin{equation}
        \sum_{k \in K_f} y_{0k} \leq |K_f|, \quad \forall f \in F \label{eq:fleet}
    \end{equation}
    约束(\ref{eq:fleet})确保每种车型的使用数量不超过可用车辆数。
\end{enumerate}

\subsubsection{求解算法}

【算法选择理由】由于问题规模较大（【具体规模】），精确算法难以在竞赛时间内求解。本文采用【算法名称】进行求解。

【算法框架描述】：
\begin{algorithm}[H]
\caption{【算法名称】主流程}
\label{alg:main}
\begin{algorithmic}[1]
\REQUIRE 【输入：初始参数、数据】
\ENSURE 【输出：最优解】
\STATE 生成初始解（【策略描述】）
\STATE 初始化温度 $T \leftarrow T_0$，迭代计数器 $iter \leftarrow 0$
\FOR{$iter = 1$ \TO $max\_iterations$}
    \STATE 自适应选择【破坏/变异/邻域】算子
    \STATE 执行【破坏/变异】操作，产生中间解
    \STATE 执行【修复/交叉】操作，生成新解 $S_{\text{new}}$
    \STATE 计算成本增量 $\Delta = \text{cost}(S_{\text{new}}) - \text{cost}(S_{\text{current}})$
    \IF{$\Delta < 0$ \OR $\text{random}() < \exp(-\Delta / T)$}
        \STATE $S_{\text{current}} \leftarrow S_{\text{new}}$
        \IF{$\text{cost}(S_{\text{new}}) < \text{cost}(S_{\text{best}})$}
            \STATE $S_{\text{best}} \leftarrow S_{\text{new}}$
        \ENDIF
    \ENDIF
    \STATE 更新算子得分与权重
    \STATE 降温 $T \leftarrow T \times \alpha$
\ENDFOR
\RETURN $S_{\text{best}}$
\end{algorithmic}
\end{algorithm}

【算法关键设计】：
\begin{enumerate}
    \item \textbf{初始解生成}：【描述策略】
    \item \textbf{邻域结构}：【描述破坏/修复算子的设计和数量】
    \item \textbf{约束处理}：【描述如何处理关键约束】
    \item \textbf{参数设置}：【列出关键参数及其取值】
\end{enumerate}

\subsubsection{求解结果}

【在此展示模型求解结果，包括主要数值指标、路线统计、成本构成等。用表格和图表辅助说明。约600-800字。】

% --- 示例：结果表格 ---
%\begin{table}[H]
%\centering
%\caption{问题1求解结果汇总}
%\label{tab:q1-results}
%\begin{tabular}{l|c|c}
%\toprule
%\textbf{指标} & \textbf{数值} & \textbf{单位} \\
%\midrule
%使用车辆数 & 【】 & 辆 \\
%运营总成本 & 【】 & 元 \\
%能耗成本 & 【】 & 元 \\
%碳排放总量 & 【】 & kg \\
%\bottomrule
%\end{tabular}
%\end{table}

% --- 5.3 问题2 ---
\subsection{问题2：【子标题】}

\subsubsection{模型建立}
【说明问题2与问题1的差异，新增的约束和修改的模型。约300-400字。】

\subsubsection{求解方法}
【说明求解策略的调整。约300-400字。】

% --- 5.4 问题3 ---
\subsection{问题3：【子标题】}

\subsubsection{【子节标题】}
【约300-400字。】

\subsubsection{求解结果}
【约300-400字。】

% ==========================================
% 六、模型的评价与优化
% ==========================================
\section{模型的评价与优化}

\subsection{验证方法}
【说明如何验证模型的正确性和结果的可信度。如：小规模实例的精确求解对比、与已知最优解的比较、约束满足度检查等。约200-300字。】

\subsection{敏感性分析}
【分析关键参数（如碳税税率、速度、时间窗宽度等）变化对结果的影响。使用图表展示灵敏度。约400-500字。】

% --- 示例：灵敏度分析 ---
%\begin{figure}[H]
%    \centering
%    \includegraphics[width=0.8\textwidth]{figure_sensitivity.png}
%    \caption{关键参数灵敏度分析}
%    \label{fig:sensitivity}
%\end{figure}

\subsection{模型优点}
\begin{enumerate}
    \item 【优点1，需具体说明而不能泛泛而谈。如"综合考虑了车辆异构性、时变速度和软时间窗，模型具有较强的现实适用性"】
    \item 【优点2，需配以具体数据或对比】
    \item 【优点3】
    \item 【优点4】
\end{enumerate}

\subsection{模型缺点}
\begin{enumerate}
    \item 【缺点1，诚实指出局限性，如"假设速度分布为已知的静态分布，未考虑实时交通状况的动态变化"】
    \item 【缺点2】
    \item 【缺点3】
\end{enumerate}

% ==========================================
% 七、模型的应用与推广
% ==========================================
\section{模型的应用与推广}

【约300-400字，讨论以下内容：】
\begin{enumerate}
    \item 模型在其他物流配送场景中的应用（如同城快递、冷链配送、应急物资调度等）
    \item 模型在其他行业的借鉴意义（如公共交通调度、共享汽车调度等）
    \item 需要做的调整和扩展（如加入实时交通数据、多配送中心、动态定价等）
\end{enumerate}

% ==========================================
% 参考文献
% ==========================================
\newpage
\begin{thebibliography}{99}

\bibitem{ref1} 作者1, 作者2. 文章标题[J]. 期刊名, 年份, 卷(期): 页码.

\bibitem{ref2} 作者. 文章标题[J]. 期刊名, 年份, 卷(期): 页码.

\bibitem{ref3} 作者. 书名[M]. 出版地: 出版社, 年份.

\bibitem{ref4} 作者. 文章标题[C]// 会议名称. 会议地点, 年份: 页码.

\bibitem{ref5} 作者. 文章标题[J]. 期刊名, 年份, 卷(期): 页码.

\end{thebibliography}

% ==========================================
% 附录
% ==========================================
\newpage
\section*{附录}
\addcontentsline{toc}{section}{附录}

% --- 附录A ---
\subsection*{附录A：【附录标题】}
\addcontentsline{toc}{subsection}{附录A：【附录标题】}

【附录正文，可包含数据表格、补充实验结果等。】

% --- 附录B ---
\subsection*{附录B：【附录标题】}
\addcontentsline{toc}{subsection}{附录B：【附录标题】}

【附录正文。】

% --- 附录C：核心求解代码 ---
\subsection*{附录C：核心求解代码}
\addcontentsline{toc}{subsection}{附录C：核心求解代码}

% 代码使用 small 字号和等宽字体
{\small\ttfamily
\begin{verbatim}
# ===== 在此处粘贴核心Python代码 =====
# 建议包含以下关键函数：
# 1. 主求解入口
# 2. 关键算法步骤
# 3. 约束处理函数
#
# 示例：
def solve(self, ...):
    """主求解入口"""
    # 初始化
    # 迭代求解
    # 返回最优解
    pass
\end{verbatim}
}

\end{document}
